\section{Tema ampliado: multitarea y deriva de distribución}
\subsection{Policy gradient multitarea}
En multi-task RL, el policy gradient necesita compartir parámetros entre distintas tareas. El enfoque más común es usar un shared encoder + task-specific head, lo que permite cambiar rápidamente entre distintas reward function. Durante el entrenamiento suele agregarse un task-id embedding para distinguir los contexts.

\begin{knowledgebox}{Puntos clave del entrenamiento multitarea}
\begin{itemize}
    \item En un mismo lote se mezclan varias tareas; si el reward scale difiere, es necesario estandarizarlo;
    \item Se usa un task embedding para que la policy sea aware of active task;
    \item El critic puede compartir el backbone, pero un value head de múltiples cabezas se usa para estimar el advantage por separado.
\end{itemize}
\end{knowledgebox}

\subsection{Deriva de distribución y domain adaptation}
En el despliegue real, la diferencia entre los datos de entrenamiento y los datos en producción provoca policy drift. Las formas comunes de mitigarlo incluyen: agregar domain randomization, usar un adversarial buffer para muestrear hard cases, y ampliar automáticamente el replay buffer durante el rollout (cohort of policies). Necesitamos registrar explícitamente en el log el `data shift ratio` y el `reward shift ratio`.

\begin{warningbox}{Consecuencias de ignorar el domain shift}
Ignorar la deriva de distribución hace que el desempeño de la política colapse ante cambios pequeños (como sensor noise), y quedar registrado en una compliance review puede clasificarse como «lacks robustness».
\end{warningbox}

\subsection{Resumen de la sección}
El entrenamiento multitarea y la deriva de distribución son los retos clave tras extender el gradiente de política a sistemas reales, y exigen trabajar a fondo tanto en la policy architecture como en el logging pipeline.

